% GENERATED by tools/gen.py from the unit's vocabulary -- do not edit.
Ćwiczenia tego bloku są \emph{generowane} z listy słów tygodnia i sprawdzają pamięć słownictwa, nie gramatykę. Jeśli odpowiedź jest synonimem, a nie słowem z listy, sprawdź w tabeli powyżej.
\begin{cwiczeniadalej}
\exitem{Podaj rodzajnik: \dots{} \de{Redewendung} (idiom, wyrażenie)}{\artf{die} Redewendung}
\exitem{Podaj rodzajnik: \dots{} \de{Sprichwort} (przysłowie)}{\artn{das} Sprichwort}
\exitem{Podaj rodzajnik: \dots{} \de{Wortbildung} (słowotwórstwo)}{\artf{die} Wortbildung}
\exitem{Podaj rodzajnik: \dots{} \de{Präfix} (przedrostek)}{\artn{das} Präfix}
\exitem{Podaj rodzajnik: \dots{} \de{Suffix} (przyrostek)}{\artn{das} Suffix}
\exitem{Podaj rodzajnik: \dots{} \de{Kompositum} (wyraz złożony)}{\artn{das} Kompositum}
\exitem{Podaj rodzajnik: \dots{} \de{Partikel} (partykuła)}{\artf{die} Partikel}
\exitem{Podaj rodzajnik: \dots{} \de{Nuance} (niuans)}{\artf{die} Nuance}
\exitem{Podaj rodzajnik: \dots{} \de{Betonung} (akcent, nacisk)}{\artf{die} Betonung}
\exitem{Podaj rodzajnik: \dots{} \de{Umgangssprache} (język potoczny)}{\artf{die} Umgangssprache}
\exitem{Podaj liczbę mnogą: \de{\artf{die} Redewendung}}{\de{\artf{die} Redewendungen}}
\exitem{Podaj liczbę mnogą: \de{\artf{die} Wortbildung}}{\de{\artf{die} Wortbildungen}}
\exitem{Podaj liczbę mnogą: \de{\artn{das} Präfix}}{\de{\artf{die} Präfixe}}
\exitem{Podaj liczbę mnogą: \de{\artn{das} Kompositum}}{\de{\artf{die} Komposita}}
\exitem{Podaj liczbę mnogą: \de{\artf{die} Partikel}}{\de{\artf{die} Partikeln}}
\exitem{Podaj liczbę mnogą: \de{\artn{das} Register}}{\de{\artf{die} Register}}
\exitem{Podaj liczbę mnogą: \de{\artf{die} Nuance}}{\de{\artf{die} Nuancen}}
\exitem{Podaj liczbę mnogą: \de{\artf{die} Betonung}}{\de{\artf{die} Betonungen}}
\exitem{Jak to jest po niemiecku: idiom, wyrażenie?}{\de{\artf{die} Redewendung, Redewendungen}}
\exitem{Jak to jest po niemiecku: słowotwórstwo?}{\de{\artf{die} Wortbildung, Wortbildungen}}
\exitem{Jak to jest po niemiecku: przedrostek?}{\de{\artn{das} Präfix, Präfixe}}
\exitem{Jak to jest po niemiecku: wyraz złożony?}{\de{\artn{das} Kompositum, Komposita}}
\exitem{Jak to jest po niemiecku: ton, intonacja?}{\de{\artm{der} Tonfall}}
\exitem{Jak to jest po niemiecku: akcent, nacisk?}{\de{\artf{die} Betonung, Betonungen}}
\exitem{Jak to jest po niemiecku: no, raz?}{\de{mal}}
\exitem{Jak to jest po niemiecku: właśnie, po prostu?}{\de{eben}}
\exitem{Jak to jest po niemiecku: po prostu?}{\de{halt}}
\exitem{Jak to jest po niemiecku: tylko, no?}{\de{bloß}}
\exitem{Co znaczy \de{\artn{das} Sprichwort, Sprichwörter}?}{przysłowie}
\exitem{Co znaczy \de{\artf{die} Wortbildung, Wortbildungen}?}{słowotwórstwo}
\exitem{Co znaczy \de{\artn{das} Präfix, Präfixe}?}{przedrostek}
\exitem{Co znaczy \de{\artn{das} Kompositum, Komposita}?}{wyraz złożony}
\exitem{Co znaczy \de{\artf{die} Partikel, Partikeln}?}{partykuła}
\exitem{Co znaczy \de{\artn{das} Register, Register}?}{rejestr (językowy)}
\exitem{Co znaczy \de{doch}?}{przecież, jednak}
\exitem{Co znaczy \de{eben}?}{właśnie, po prostu}
\exitem{Co znaczy \de{halt}?}{po prostu}
\exitem{Co znaczy \de{nämlich}?}{mianowicie}
\end{cwiczeniadalej}
